

Eucalyptus is a tree that, while nowadays common in the Mediterranean landscape, has only been around Europe for less than 200 years. Originally from Australia, it was first discovered by Europeans in the late 18th century, following the maritime expeditions of James Cook, through which some samples of the plant were brought over to Europe. It was based on these that the botanist Charles Louis L'Héritier gave the plant its current name: "Eucalyptus". However, it is not until the second half of the 19th century that the tree definitely entered the Mediterranean landscape, notably through the works of Prosper Ramel, who, after being key in the import of the tree to Europe, started large plantations of eucalyptus in France and Algeria\footnote{\cite[2]{davin-mortier_leucalyptus_2026}}. These new plantations, in turn, fuel a discourse about this new exotic tree, which had many virtues and qualities; both industrial (strong wood for construction) and domestic (folk medicine) in nature.

%TODO: CHANGE CITATION AND ADD THE GALICIAN BOOK JUST IN CASE: Pablo Ramil Rego, ‘O Eucalipto En Galicia’, Recursos Rurais: Revista Oficial Do Instituto de Biodiversidade Agraria e Desenvolvemento Rural (IBADER), no. 15 (2019), 5.

Current research is lacking in the analysis of the discourse around the importation and the spread of eucalyptus throughout the French-speaking world. We can cite the work of Robin W. Doughty\footnote{\cite{doughty_eucalyptus_2000}}, tracing the history of eucalyptus and its usages, however more from an anglophone perspective. We can also cite the work of Elisabeth Davin-Mortier, who focuses more on Europe, the Middle-East and more generally on the Mediterranean basin, looking at the links between the plant and its role in European colonization. Thus, the aim of this project is to address this gap, via a quantitative approach, by analyzing more than 17'000 different documents, published between 1783 and 1920. We chose to limit the corpus at the end of the first world war, as it represents a key event impacting the regions studied, at a time where we can already find large eucalyptus plantations around the Mediterranean. So, in this project, we will try to understand how and through which actors the discourse around eucalyptus is shaped. More precisely, understand in which contexts the tree is mentioned and with what kind of properties it is associated in the documents. The analysis will be pursued both through time and space, to see how the discourse evolves, and from where it originates. 
